\section{Differentialrechnung}


\Title{Differenzierbarkeit}

\begin{tcolorbox}
    Sei $x_0 \in D$ ein Häufungspunkt, so $f$ ist in $x_0$ \textbf{differenzierbar}, falls der Grenzwert
    $$\Lim_{x \to x_0} \frac{f(x) - f(x_0)}{x - x_0} = \Lim_{h \to 0} \frac{f(x_0 + h) - f(x_0)}{h}$$
    existiert. Der Grenzwert wird dann mit $f'(x_0)$ oder $\frac{df}{dx}(x_0)$ bezeichnet.
\end{tcolorbox}

\begin{tcolorbox}
	\Satz (Weierstrass) Sei $f: D \rightarrow R$, $x_{0} \in D$ Häufungspunkt von $D$. Dann ist folgende
    Aussage äquivalent zur differenzierbarkeit:

	$\exists c \in \R$ und $r: D \rightarrow \R$ mit
	\begin{itemize}
		\item{$f(x) = f(x_0) + c(x-x_0) + r(x)(x-x_0)$}
		\item{$r(x_0)=0$ und r ist stetig in $x_0$}
    \end{itemize}
	Falls dies zutrifft ist $c=f^{'}(x_0)$ eindeutig bestimmt. \bigskip

	\Satz Eine Funktion $f: D \rightarrow \R$ ist genau dann in $x_0$ differenzierbar falls $\exists$ 
    Funktion $\phi(x) := f'(x_0) + r(x)$, so dass $f(x) = f(x_0) + \phi(x)(x-x_0), \; \forall x \in D$. 
    In diesem Fall gilt $\phi(x_0) = f^{'}(x_0)$. Wenn man die Gleichung umformt erhält man $\phi(x)
    = \frac{f(x) - f(x_0)}{x - x_0}$.
\end{tcolorbox}

\Kor $f$ differenzierbar in $x_0 \implies$ $f$ stetig in $x_0$ 

\Bem Die Tangentengleichung von $f$ im Punkt $(x_0, f(x_0))$ ist $y = f(x_0) + f'(x_0)(x - x_0)$.

\begin{tcolorbox}
    $f : D \to R$ ist in \textbf{$\mathbf{D}$ differenzierbar}, falls für jeden Häufungspunkt $x_0 \in D$,
    $f$ in $x_0$ differenzierbar ist.
\end{tcolorbox}

\Kor Sei $f: D \to E$ eine bijektive Funktion und $x_0 \in D$ ein Häufungspunkt. Wir nehmen an $f$ ist
in $x_0$ differenzierbar und $f'(x_0) \neq 0$, zudem nehmen wir an $f^{-1}$ ist in $y_0 = f(x_0)$
stetig. Dann ist $y_0$ ein Häufungspunkt von $E$ und $f^{-1}$ ist in $y_0$ differenzierbar.
$$(f^{-1})'(y_0) = \frac{1}{f'(f^{-1}(y_0))}$$

\Title{Rechnen mit Ableitungen}

\begin{enumerate}[label = (\arabic*)]
    \item $(f + g)'(x_0) = f'(x_0) + g'(x_0)$
    \item $(f \cdot g)'(x_0) = f'(x_0) \cdot g(x_0) + g'(x_0) \cdot f(x_0)$
    \item $(\frac{f}{g})'(x_0) = \frac{f'(x_0)g(x_0) - f(x_0)g'(x_0)}{g(x_0)^2}$ für $g(x_0) \neq 0$
    \item $(g \circ f)'(x_0) = g'(f(x_0)) \cdot f'(x_0)$
    \item $(f^{-1})'(y_0) = \frac{1}{f'(x_0)}$, wobei $y_0 = f(x_0)$ und $f'(x_0) \neq 0$
    \item $(a^{f(x)})' = ln(a)\cdot a^{f(x)} \cdot f`(x)$
    \item $(f(x)^{g(x)})' = f(x)^{g(x)} \cdot [ln(f(x))\cdot g(x)]'$
\end{enumerate}

\Bem Für gerade $k$ gilt $\cosh(x)^{(k)} = \cosh(x)$ und für ungerade $k$ gilt $\cosh(x)^{(k)} = \sinh(x)$,
analoges gilt für $\sinh$.


\columnbreak
%----------------------------------------------------------------------------------------------------


\Title{Aussagen der Ableitung}

Wir nehmen an $f$ ist in $D$ differenzierbar, so können wir folgende Aussagen machen:
\begin{enumerate}[label = (\arabic*)]
    \item $f$ besitzt ein lokales Maximum in $x_0$, wenn $f'(x_0) = 0$ und $f''(x_0) < 0$ oder falls es $\delta > 0$ gibt mit:
    $$f(x) \leq f(x_0) \; \; \forall x \in ]x_0 - \delta, x_0 + \delta[ \cap D$$
    \item $f$ besitzt ein lokales Minimum in $x_0$, wenn $f'(x_0) = 0$ und $f''(x_0) > 0$ oder falls es $\delta > 0$ gibt mit:
    $$f(x) \geq f(x_0) \; \; \forall x \in ]x_0 - \delta, x_0 + \delta[ \cap D$$
    \item $f$ besitzt ein lokales Extremum in $x_0$ falls es entweder ein lokales Maximum oder Minimum
    von $f$ ist. Dies ist der Fall wenn $f'(x_0) = 0$ und $f''(x_0) \neq 0$.
    \item $f$ besitzt einen Wendepunkt wenn $f''(x_0) = 0$.
    \item $f$ besitzt einen Sattelpunkt wenn $f'(x_0) = 0$ und $f''(x_0) = 0$.
\end{enumerate}

Weiter gelten auch folgende Aussagen für $f,g : [a,b] \to \R$, wobei $f,g$ stetig und in $]a,b[$
differenzierbar sind.
\begin{enumerate}[label = (\arabic*)]
    \item Falls $f'(\epsilon) = 0 \; \; \forall \epsilon \in ]a,b[$ ist $f$ konstant
    \item Falls $f'(\epsilon) = g'(\epsilon) \; \; \forall \epsilon \in ]a,b[$ ist $f = g + c$
    \item Falls $f'(\epsilon) \leq [<] 0 \; \; \forall \epsilon \in ]a,b[$ ist $f$ [streng] monoton fallend
    \item Falls $f'(\epsilon) \geq [>] 0 \; \; \forall \epsilon \in ]a,b[$ ist $f$ [streng] monoton wachsend
\end{enumerate}


\Title{Satz von Rolle}

\Satz Sei $f : [a,b] \to \R$ stetig und in $]a,b[$ differenzierbar. Wenn nun $f(a) = f(b)$ gilt, so 
gibt es ein $\epsilon \in ]a,b[$ mit $f'(\epsilon) = 0$.


\Title{Mittelwertsatz}

\Satz Sei $f : [a,b] \to \R$ stetig und in $]a,b[$ differenzierbar. Dann gibt es $\epsilon \in ]a,b[$ 
mit $f(b) - f(a) = f'(\epsilon)(b - a)$.


\Title{Satz von L'Hospital}

Sei $f : [a,b] \to \R$ stetig und in $]a,b[$ differenzierbar. Dann gibt es $\epsilon \in ]a,b[$ mit
$g'(\epsilon)(f(b)-f(a)) = f'(\epsilon)(g(b)-g(a))$. Falls nun $g'(x) \neq 0$ für alle $x$ ist, dann
folgt $g(a) \neq g(b)$ und 
$$\frac{f(b) - f(a)}{g(b) - g(a)} = \frac{f'(\epsilon)}{g'(\epsilon)}.$$

\Satz Falls nun 
$$\Lim_{x \to b^-} f(x) = 0, \Lim_{x \to b^-} g(x) = 0 \; \text{ und } \; \Lim_{x \to b^-} \frac{f'(x)}{g'(x)} = \lambda$$
existiert, folgt 
$$\Lim_{x \to b^-} \frac{f(x)}{g(x)} = \Lim_{x \to b^-} \frac{f'(x)}{g'(x)} = \lambda $$ 

L'Hospital gilt auch wenn: $b = +\infty$, $x \to a^+$, $\lambda = +\infty$ oder 
$\lim f = \lim g = \infty$.

Weiter gilt die Regel von L'Hospital auch für Fälle der Form ``$\frac{\infty}{\infty}$'', indem man
sie wiefolgt umschreibt: $\frac{f(x)}{g(x)}=\frac 1{g(x)} / \frac 1{f(x)}$.
Für Fälle der Form ``$\infty-\infty$'' kann man die Funktionen auf den gleichen Nenner bringen.

\Bsp $\lim_{x \to 0}\left(\frac1{\sin x}-\frac1x\right)=\lim_{x \to 0}\frac{x-\sin x}{x\sin x}$


\columnbreak
%----------------------------------------------------------------------------------------------------


\Title{Konvexität}

$f$ ist [streng] \textbf{konvex} (auf $I$) falls für alle $x \leq [<] y$, $x,y \in I$ und $\lambda \in [0,1]$
$$f(\lambda x + (1-\lambda)y) \quad \leq [<] \quad \lambda f(x) + (1-\lambda)f(y)$$
gilt.

Wobei wir folgende Bemerkungen machen können:
\begin{enumerate}[label = (\arabic*)]
    \item Die Summe zweier konvexer Funktionen ist konvex
    \item $f$ ist genau dann konvex, falls für $x_0 \leq x \leq x_1$ in $I$ gilt:
    $$\frac{f(x) - f(x_0)}{x - x_0} \leq \frac{f(x_1) - f(x)}{x_1 - x}$$
    \item Die Funktion $f$ ist genau dann [streng] konvex, falls $f'$ [streng] monoton wachsend ist.
\end{enumerate}

\Bem Alle Aussagen gelten auch für \textbf{Konkavität}, wir müssen nur das Ungleichzeichen umkehren.


\Title{Höhere Ableitungen}

Sei $f : D \to \R$ differenzierbar.
\begin{enumerate}[label = (\arabic*)]
    \item Für $n \geq 2$ ist $f$ \textbf{$n$-mal differenzierbar} in $D$, falls $f^{(n-1)}$ in $D$ 
    differenzierbar ist. Wobei $n$-mal differenzierbar $\implies$ $(n-1)$-mal stetig differenzierbar.
    \item Die Funktion $f$ ist in $D$ \textbf{glatt}, falls sie $\forall n \geq 1$, $n$-mal 
    differenzierbar ist.
\end{enumerate}

\Bem Alle Polynome, sowie $e^x, \sin$ und $\cos$ sind glatte Funktionen.


\Title{Rechnen mit Höhere Ableitungen}

Seien $f,g : D \to \R$ $n$-mal differenzierbar:
\begin{enumerate}[label = (\arabic*)]
    \item $(f + g)^{(n)} = f^{(n)} + g^{(n)}$
    \item $(f \cdot g)^{(n)} = \Sum_{k = 0}^{n} \binom{n}{k} f^{(k)} \cdot g^{(n-k)}$
    \item $\frac{f}{g}$ ist $n$-mal differenzierbar falls $g(x) \neq 0, \forall x \in D$
    \item $(g \circ f)$ ist $n$-mal differenzierbar
\end{enumerate}


\Title{Potenzreihe}

\Satz Sei $f_n$ eine Funktionenfolge, wobei $f_n$ einmal stetig differenzierbar ist für alle 
$n \geq 1$. Wir nehmen an, dass sowohl die Folge $(f_n)_{n \geq 1}$ wie $(f_n')_{n \geq 1}$ 
gleichmässig konvergieren mit $\Limit f_n = f$ und $\Limit f_n' = p$. Dann ist $f$ stetig 
differenzierbar und $f' = p$.

\begin{tcolorbox}
    \Satz Sei $\Reihe c_k x^k$ eine Potenzreihe mit Konvergenzradius $p > 0$. Dann ist 
    $$f(x) = \Reihe c_k (x - x_0)^k$$
    auf $]x_0 - p, x_0 + p[$ differenzierbar und
    $$f'(x) = \Reihe k \cdot c_k (x - x_0)^{k-1}$$
    für alle $x \in ]x_0 - p, x_0 + p[$.
\end{tcolorbox}


\columnbreak
%----------------------------------------------------------------------------------------------------


\Title{Taylox Approximation}

Jede glatte Funktione kann als Potenzreihe angenähert werden, dafür brauchen wir ein sogenanntes
\textbf{Taylor-Polynom} (oder Taylorreihen). Die trigonometrischen Funktionen sind genau solche 
Taylorreihen.

\begin{tcolorbox}
    Sei $f : [a,b] \to \R$ stetig und in $]a,b[$ $(n+1)$-mal differenzierbar. Für jedes $a < x \leq b$
    gibt es $\epsilon \in ]a,x[$ mit:
    $$f(x) = \Sum_{k = 0}^{n} \frac{f^{(k)}(a)}{k!}(x - a)^k + \frac{f^{(n + 1)}(\epsilon)}{(n+1)!}(x-a)^{n+1}$$
    Wobei wir hierbei das $n$-te Taylor-Polynom $T_n$ und eine Fehlerabschätzung $R_n$ erhalten.
\end{tcolorbox}

\Bem Der letzte Term $\frac{f^{(n + 1)}(\epsilon)}{(n+1)!}(x-a)^{n+1}$ wird meist zur Fehlerabschätzung
innerhalb eines Bereiches von $a$ verwendet. Man nimmt daher den Wert von $\epsilon \in ]a,x[$, so
dass der Fehlerterm am grössten ist.


\Title{Spezielle Punkte}

Sei $n \geq 0, a < x_0 < b$ und $f : [a,b] \to \R$ in $]a,b[$ $(n+1)$-mal stetig differenzierbar.
Annahme: $f'(x_0) = f''(x_0) = ... = f^{(n)}(x_0) = 0$

\begin{enumerate}[label = (\arabic*)]
    \item Falls $n$ gerade ist und $x_0$ eine lokale Extremalstelle, folgt $f^{(n+1)}(x_0) = 0$
    \item Falls $n$ ungerade ist und $f^{(n+1)}(x_0) > 0$, so ist $x_0$ eine strikte lokale Minimalstelle
    \item Falls $n$ ungerade ist und $f^{(n+1)}(x_0) < 0$, so ist $x_0$ eine strikte lokale Maximalstelle
\end{enumerate}

Ist $x_0$ jedoch keine Extremalstelle bleiben zwei Optionen:
\begin{enumerate}[label = (\arabic*)]
    \item $f'(x_0) = 0 \wedge x_0$ keine Extremalstelle $\implies$ ist ein Sattelpunkt
    \item $f''(x_0) = 0 \wedge x_0$ keine Extremalstelle $\implies$ ist ein Wendepunkt
\end{enumerate}
